\documentclass{article}

\usepackage[utf8]{inputenc}
\usepackage[russian]{babel} % For Russian language support
\usepackage{amsmath}
\usepackage{amsfonts}
\usepackage{amssymb}

\title{Пример научной статьи}
\author{Автор}
\date{\today}

\begin{document}

\maketitle

\section{Введение}
Научные статьи играют важную роль в распространении знаний и результатов исследований. 
Введение обычно содержит основную информацию о теме статьи, целях исследования и его актуальности. 
В этой статье мы рассмотрим структуру научной статьи и основные компоненты, которые должны в ней присутствовать.

\section{Основная часть}
Основная часть научной статьи содержит детальное описание методов исследования, 
результатов экспериментов и анализа данных. В этой части авторы обычно представляют 
свою точку зрения, аргументируют свою позицию и ссылаются на предыдущие работы в этой области.

\section{Заключение}
Заключение подводит итоги всей работы, обобщает полученные результаты и 
определяет направления для дальнейших исследований. В заключении также 
можно указать практическое применение полученных результатов и их значение 
для науки и общества. Таким образом, научная статья должна иметь четкую 
структуру, логичное изложение материала и содержательные выводы.

\end{document}